\documentclass[11pt, oneside]{article}
\usepackage{geometry} 
\geometry{letterpaper}
\usepackage{graphicx}	

\usepackage{amssymb, amsmath}

\newcommand{\I}[2]{
  \ensuremath{ \mathbb{I}_{#1} \! \left[ #2 \right] }
}

\newcommand{\E}[2]{
  \ensuremath{ \mathbb{E}_{#1} \! \left[ #2 \right] }
}

\newcommand{\rnd}[2]{
  \ensuremath{ \frac{ \mathrm{d} #1}{ \mathrm{d} #2 } }
}

\begin{document}

\section{Integral Notation}

Partially evaluating a conditional probability distribution on any $y \in Y$ 
yields a probability distribution that concentrates on the level set 
$f^{-1}(y)$,
%
\begin{alignat*}{6}
\pi^{f}_{y} :\; &\mathcal{X}& &\rightarrow& \; &[0, 1]&
\\
&\mathsf{x}& &\mapsto& &\pi^{f}( \mathsf{x} \mid y )&.
\end{alignat*}

When paired with an integrand $g : X \rightarrow \mathbb{R}$ this collection 
of probability distributions defines a \textbf{conditional expectation function},
%
\begin{alignat*}{6}
e_{g} : \; &Y& &\rightarrow& \; &\mathbb{R}&
\\
&y& &\mapsto& & \E{ \pi^{f}_{y} }{ g }&.
\end{alignat*}

The law of total expectation ensures that the pushforward expectation of this 
conditional expectation function is always equal to the original expectation 
value,
%
\begin{equation*}
\E{\pi}{g} = \E{ f_{*} \pi }{e_{g}}.
\end{equation*}

Ideally we'd be able to write the law of total expectation more compactly,
packing everything into one line.  Unfortunately any compact notation is 
frustrated by the standard notational convention of ignoring arguments when 
denoting expectation values.  For example a notation that replaces $y$ with 
a placeholder ``$\cdot$",
%
\begin{equation*}
\E{\pi}{g} = \E{ f_{*} \pi }{e_{g}} = \E{ f_{*} \pi }{ \E{\pi^{f}_{\cdot}}{g} },
\end{equation*}
%
is not only hard to read but can be ambiguous in applications where there 
are multiple spaces on which $f$ and $g$ might act.

At the same time a more explicit notation like
%
\begin{equation*}
\E{\pi}{g} = \E{ f_{*} \pi }{e_{g}} = \E{ f_{*} \pi }{ \E{\pi^{f}_{y}}{g(x)} }
\end{equation*}
%
breaks the standard notational convention without clarifying on what spaces 
the probability distributions $f_{*} \pi$ and $\pi^{f}_{y}$ are defined.

Conditional probability theory is one place where the more explicit integral 
notation for expectation values that we first discussed in \textbf{previous chapter}
becomes incredibly useful.  To make the notation as compact as possible we
will need to  introduce a \emph{conditional variable} $z_{y}$ that takes values 
in only the level set corresponding to the output point $y \in Y$,
%
\begin{equation*}
z_{y} \in f^{-1}(y) \subset X.
\end{equation*}
%
Conditional variables allow us to decompose any input variable into an output 
variable and a corresponding conditional variable,
%
\begin{equation*}
x = (y, z_{y}).
\end{equation*}
%
Note that the right hand-side isn't an ordered pair because the possible
values of the second variable will in general depend on the choice of the 
first variable.  For the mathematically-curious this construction is known 
as a \textbf{semi-direct product}.

Using conditional variables we can write the conditional expectation value as
%
\begin{align*}
e_{g}(y) 
&= 
\E{ \pi^{f}_{y} }{ g }
\\
&=
\int \pi^{f}( \mathrm{d} x \mid y ) \, g(x),
\end{align*}
%
with the law of total expectation becoming
%
\begin{align*}
\E{\pi}{g} 
&= 
\E{ f_{*} \pi }{e_{g}}
\\
\int \pi( \mathrm{d} x ) \, g(x)
&=
\int f_{*} \pi (\mathrm{d} y) \, e_{g}(y)
\\
\int \pi( \mathrm{d} x ) \, g(x)
&=
\int f_{*} \pi (\mathrm{d} y) \int \pi^{f}( \mathrm{d} x \mid y ) \, g(x).
\end{align*}
%
\textbf{Convert $\mathrm{d} x \mid y$ to $\mathrm{d} x_{y} \mid y$.}

I will make extensive use of this compact notation anytime conditional 
probability theory might come into play which, as we will see, will be 
often.

\section{Conditional Probability Density Functions}

\subsection{Non-Null Partitions}

Constructing conditional probability density functions for non-null 
partitions is relatively straightforward.  Recall that in this case 
we can define a conditional probability as
%
\begin{equation*}
\pi^{\mathcal{P}}( \mathsf{x} \mid \mathsf{c}_{j} )
=
\frac{ \pi( \mathsf{x} \cap \mathsf{c}_{j} ) }{ \pi( \mathsf{c}_{j} ) }.
\end{equation*}
%
More generally for any integrand $g : X \rightarrow \mathbb{R}$ these 
conditional probability allocations define a conditional expectation 
value,
%
\begin{equation*}
\E{ \pi^{\mathcal{P}}_{\mathsf{c}_{j}} }{g}
=
\frac{ \E{\pi}{ I_{\mathsf{c}_{j}} \cdot g } }
{ \E{\pi}{ I_{\mathsf{c}_{j}} } }.
\end{equation*}

Given a convenient reference measure $\mu$ we can convert these expectation 
values into $\mu$-informed integrals with the introduction of a Radon-Nikodym
derivative.  On the left-hand side we have
%
\begin{equation*}
\E{\pi^{\mathcal{P}}_{\mathsf{c}_{j}}}{g}
=
\I{\mu}{ \rnd{\pi^{\mathcal{P}}_{\mathsf{c}_{j}}}{\mu} \cdot g }
\end{equation*}
%
and on the right-hand side we have
%
\begin{align*}
\frac{ \E{\pi}{ I_{\mathsf{c}_{j}} \cdot g } }
{ \E{\pi}{ I_{\mathsf{c}_{j}} } }
&=
\frac{ \I{\mu}{ \rnd{\pi}{\mu} \cdot I_{\mathsf{c}_{j}} \cdot g } }
{ \I{\mu}{ \rnd{\pi}{\mu} \cdot I_{\mathsf{c}_{j}} } }
\\
&=
\I{\mu}{ 
\frac{ \rnd{\pi}{\mu} \cdot I_{\mathsf{c}_{j}} }
{ \I{\mu}{ \rnd{\pi}{\mu} \cdot I_{\mathsf{c}_{j}} } 
} \cdot g }.
\end{align*}

Consequently
%
\begin{align*}
\E{\pi^{\mathcal{P}}_{\mathsf{c}_{j}}}{g}
&=
\frac{ \E{\pi}{ I_{\mathsf{c}_{j}} \cdot g } }
{ \E{\pi}{ I_{\mathsf{c}_{j}} } }
\\
\I{\mu}{ \rnd{\pi^{\mathcal{P}}_{\mathsf{c}_{j}}}{\mu} \cdot g }
&=
\I{\mu}{ 
\frac{ \rnd{\pi}{\mu} \cdot I_{\mathsf{c}_{j}} }
{ \I{\mu}{ \rnd{\pi}{\mu} \cdot I_{\mathsf{c}_{j}} } 
} \cdot g }
\end{align*}
%
which is true for any integrand $g : X \rightarrow \mathbb{R}$ if, and only if, 
%
\begin{equation*}
\rnd{\pi^{\mathcal{P}}_{\mathsf{c}_{j}}}{\mu}
\overset{\mu}{=}
\frac{\rnd{\pi}{\mu} \cdot I_{\mathsf{c}_{j}} }
{ \I{\mu}{ \rnd{\pi}{\mu} \cdot I_{\mathsf{c}_{j}} } }.
\end{equation*}
%
Repeating this calculation for each partition cell $\mathsf{c}_{j}$ defines a family 
of probability density functions that completely specifies the conditional 
probability distribution in the context of the reference measure $\mu$.

In words a conditional probability density function for a non-null partition is given 
by truncating the unconditional probability density function to each partition cell 
and then scaling by the reciprocal of the cell probability to ensure the proper 
normalization.

\textbf{figure}

\subsection{The Problem with Null Partitions}

Defining conditional probability density functions when the partition cells might 
be allocated vanishing probability is, unfortunately, a bit more complicated.
Because we can no longer define conditional probabilities as a ratio of unconditional 
probabilities we need to consult the law of total expectation in order to define 
conditional probability density functions,
%
\begin{equation*}
\int \pi( \mathrm{d} x ) \, g(x)
=
\int f_{*} \pi (\mathrm{d} y) \int \pi^{f}( \mathrm{d}x \mid y ) \, g(x),
\end{equation*}

In order to move forwards we need to convert the inner expectation into a reference 
measure-informed integral.  The subtlety, however, is exactly \emph{how} we can 
integrate over the level sets.

Consider, for example, a reference measure $\lambda$ over the entire input space
for which each conditional probability distribution $\pi^{f}( \mathsf{x} \mid y)$
is absolutely continuous.  In theory we can write the expectation over any 
particular level set as a $\lambda$-informed integral using a Radon-Nikodym
derivative,
%
\begin{equation*}
\int \pi^{f}( \mathrm{d}x \mid y ) \, g(x)
=
\int \lambda( \mathrm{d} x ) \, \rnd{ \pi^{f} }{ \lambda} (x \mid y) \, g(x).
\end{equation*}

\textbf{Actually not mathematically valid}.
While this is mathematically valid the necessary Radon-Nikodym derivative 
$\mathrm{d} \pi^{f} / \mathrm{d} \lambda$ is somewhat weird object.  For example 
remember that any subset that doesn't intersect with a given level set is allocated 
zero conditional probability,
%
\begin{align*}
0
&=
\int \pi^{f}( \mathrm{d}x \mid y ) \, I_{\mathsf{x}}(x)
\\
&=
\int \lambda( \mathrm{d} x ) \, \rnd{ \pi^{f} }{ \lambda } (x \mid y) \, I_{\mathsf{x}}(x).
\end{align*}
%
This is true, however, if and only if
%
\begin{equation*}
\rnd{ \pi^{f} }{ \lambda } (x \mid y) \overset{\lambda, \nu}{=} 0.
\end{equation*}
%
In other words $\mathrm{d} \pi^{f} / \mathrm{d} \lambda(x \mid y)$ has to be effectively 
zero outside of the level set $f^{-1}(y)$.

At the same time if the level set is allocated vanishing reference measure,
%
\begin{equation*}
\lambda( f^{-1}(y) ) = 0
\end{equation*}
%
then any integrand that is non-zero only within the level set $f^{-1})(y)$ should 
give the same integral as an integrand that is identically zero within that level 
set.  Consequently we should have
%
\begin{align*}
\int \pi^{f}( \mathrm{d}x \mid y ) \, g(x)
&=
\int \lambda( \mathrm{d} x ) \, \rnd{ \pi^{f} }{ \lambda} (x \mid y) \, g(x)
\\
&=
\int \lambda( \mathrm{d} x ) \, 0 \cdot g(x)
\\
&=
0
\end{align*}
%
regardless of $g$.  That vanishing integral then propagates through the law of 
total expectation to give
%
\begin{align*}
\int \pi( \mathrm{d} x ) \, g(x)
&=
\int f_{*} \pi (\mathrm{d} y) \int \pi^{f}( \mathrm{d}x \mid y ) \, g(x)
\\
&=
\int f_{*} \pi (\mathrm{d} y) \, 0
\\
&=
0
\end{align*}
%
for all expectands $g$, which cannot hold for arbitrary probability distributions
$\pi$.

In order to ensure non-trivial expectations the object 
$\mathrm{d} \pi^{f} / \mathrm{d} \lambda(x \mid y)$ cannot be a regular function.
Instead it would have to be a generalized function that can inflate a null subset 
to have non-vanishing integrals, mirroring the construction of the Dirac delta
``function'' that we encountered in \textbf{previous chapter}.

In practice, however, it's much more straightforward if we restrict the inner 
expectation in the law of total probability to particular level sets and define 
conditional probability density functions relative to a family of reference measures, 
each of which is restricted to a particular level set,
%
\begin{equation*}
\int \pi^{f}( \mathrm{d}x \mid y ) \, g(x)
=
\int \eta_{y}( \mathrm{d} x ) \, \rnd{ \pi^{f} }{ \eta_{y} } (x \mid y) \, g(x).
\end{equation*}

In order to relate this to probability density functions defined over the input 
and output spaces we'll need to learn how to disintegrate measures.

\subsection{Disintegrating Measures}

Disintegrating measures, however, is more subtle than disintegrating probability 
distributions.

Firstly in order to avoid the kinds of ambiguous results that arise when we allow
infinities into calculations we need to restrict our consideration to measures that 
are $\sigma$-finite.  Moreover when generalizing the law of total expectation into 
a law of total integration we cannot rely on the pushforward of an initial  
$\sigma$-finite measure because it will not always be $\sigma$-finite.

For example consider a rigid two-dimensional real space $\mathbb{R}^{2}$ equipped
with the two-dimensional Lebesgue measure $\lambda^{2}$ and a projection function
\begin{alignat*}{6}
\varpi_{1} :\; &\mathbb{R}^{2}& &\rightarrow& \; &\mathbb{R}&
\\
&(x_{1}, x_{2})& &\mapsto& &x_{1}&.
\end{alignat*}
The Lebesgue measure $\lambda_{2}$ is $\sigma$-finite, allocating finite measure 
to every measurable subset that can be encapsulated in a finite rectangle: if
%
\begin{equation*}
\mathsf{x} \subset [0, 1] \times [0, 1]
\end{equation*}
%
then
\begin{align*}
\lambda^{2}(\mathsf{x}) 
&< 
\lambda^{2}( [0, 1] \times [0, 1] )
\\
&<
l([0, 1]) \cdot l([0, 1])
\\
&<
1 \cdot 1
\\
&< 
1.
\end{align*}
Pushing $\lambda^{2}$ forward along $\varpi_{1}$, however, results in a measure 
that allocates infinite measure to finite intervals,
%
\begin{align*}
(\varpi_{1})_{*} \lambda^{2} ( [0, 1] ) 
&= 
\lambda^{2}( \varpi_{1}^{*} [0, 1]) 
\\
&= \lambda^{2}( [0, 1] \times (\infty, \infty) ) 
\\
&= l([0, 1]) \cdot l( (-\infty, \infty) )
\\
&= 1 \cdot \infty
\\
&= \infty.
\end{align*}

\textbf{figure}

In order to define a well-behaved law of total integration for $\sigma$-finite measures 
we need to specify not only a measurable function $f: X \rightarrow Y$ but also a 
$\sigma$-finite reference measure over the output space $Y$.  Formally given 
%
\begin{enumerate}
  \item an input measurable space $(X, \mathcal{X})$,
  \item a $\sigma$-finite measure $\mu : \mathcal{X} \rightarrow [0, \infty]$, 
  \item an output measurable space $(Y, \mathcal{Y})$ where every atomic subset is measurable, 
  \item a $(\mathcal{X}, \mathcal{Y})$-measurable function $f: X \rightarrow Y$, 
  \item and finally a $\sigma$-finite measure $\nu : \mathcal{Y} \rightarrow [0, \infty]$ 
\end{enumerate}
there exists at least one conditional measure
%
\begin{alignat*}{6}
\mu^{f, \nu} :\; &\mathcal{X} \times Y& &\rightarrow& \; &[0, \infty]&
\\
&\mathsf{x}, y& &\mapsto& &\mu^{f}( \mathsf{x} \mid y )&
\end{alignat*}
%
that defines a $(\mathcal{Y}, \mathcal{B}_{\mathbb{R}})$-measurable function when 
partially evaluated on the first argument,
%
\begin{alignat*}{6}
\mu^{f, \nu}_{\mathsf{x}} :\; &Y& &\rightarrow& \; &[0, \infty]&
\\
&y& &\mapsto& &\mu^{f, \nu}( \mathsf{x} \mid y )&
\end{alignat*}
%
for all $\mathsf{x} \in \mathcal{X}$, and a $\sigma$-finite measure when partially 
evaluated on the second argument,
%
\begin{alignat*}{6}
\mu^{f}_{y, \nu} :\; &\mathcal{X}& &\rightarrow& \; &[0, \infty]&
\\
&\mathsf{x}& &\mapsto& &\mu^{f, \nu}( \mathsf{x} \mid y )&
\end{alignat*}
%
$\nu$-almost all $y \in Y$.
%
Any conditional measure almost always concentrates on a given level set,
%
\begin{equation*}
\mu^{f, \nu}_{y}( \mathsf{x} ) \overset{\nu}{=} 0
\end{equation*}
%
for any $\mathsf{x} \in \mathcal{X}$ with $\mathsf{x} \cap f^{-1}(y)$
while also satisfying a law of total integration,
%
\begin{equation*}
\int \mu( \mathrm{d} x ) \, g(x)
=
\int \nu (\mathrm{d} y) \int \mu^{f, \nu}( \mathrm{d}x_{y} \mid y ) \, g(x),
\end{equation*}
%
for any well-behaved integrand $g : X \rightarrow \mathbb{R}$.

If $f_{*} \mu$ is $\sigma$-finite then we can always take $\nu = f_{*} \mu$ so that 
the law of total integration mirrors the law of total expectation.  This is always 
possible if $\mu$ is a finite measure, with $\mu(X) < \infty$, but not always
viable when $\mu$ is only $\sigma$-finite.  In particular we have to be careful 
when disintegrating Lebesgue and counting measures!



\section{Explicit Formula For Pushforward Probability Density Functions}

Disintegrating a reference measure also gives us an explicit formula for pushforward 
probability density functions.  We start with the definition of pullback expectation values:
for sufficiently measurable $f : X \rightarrow Y$ and $h : Y \rightarrow \mathbb{R}$ we have
%
\begin{equation*}
\E{\pi}{ h \circ f } = \E{f_{*} \pi}{ h },
\end{equation*}
%
or, equivalently,
%
\begin{equation*}
\int \pi( \mathrm{d} x ) \, h(f(x)) = \int f_{*} \pi( \mathrm{d} y )  \, h(y).
\end{equation*}

Our goal will be to convert these expectations into equivalent integrals so that we can 
compare the integrands, at least up to null subsets.  For example given a convenient 
reference measure $\nu$ over $Y$ we can write the right-hand side as
%
\begin{align*}
\int f_{*} \pi( \mathrm{d} y )  \, h(y)
&=
\int \nu( \mathrm{d} y ) \, \rnd{f_{*}\pi}{\nu}(y) \, h(y)
\\
&=
\int \nu( \mathrm{d} y ) \, \left[ \rnd{f_{*}\pi}{\nu}(y) \right] \, h(y).
\end{align*}

Similarly given a convenient reference measure $\mu$ over $X$ we can write the left-hand
side as
%
\begin{equation*}
\int \pi( \mathrm{d} x ) \, h(f(x))
=
\int \mu( \mathrm{d} x ) \, \rnd{\pi}{\mu}(x) \, h(f(x)).
\end{equation*}
%
Disintegrating $\mu$ with respect to $f$ and $\nu$ then gives
%
\begin{align*}
\int \pi( \mathrm{d} x ) \, h(f(x))
&=
\int \mu( \mathrm{d} x ) \, \rnd{\pi}{\mu}(x) \, h(f(x))
\\
&=
\int \nu( \mathrm{d} y ) \, \int \mu^{f, \nu}( \mathrm{d} x \mid y ) \, \rnd{\pi}{\mu}(x) \, h(f(x)).
\end{align*}
%
Because the function $h \circ f : X \rightarrow \mathbb{R}$ yields the same output for any 
$x \in f^{-1}(y)$ it is a constant with respect to the inner integral that can be factored out,
%
\begin{align*}
\int \pi( \mathrm{d} x ) \, h(f(x))
&=
\int \nu( \mathrm{d} y ) \, \int \mu^{f, \nu}( \mathrm{d} x \mid y ) \, \rnd{\pi}{\mu}(x) \, h(f(x))
\\
&=
\int \nu( \mathrm{d} y ) \, 
\left[ \int \mu^{f, \nu}( \mathrm{d} x \mid y ) \, \rnd{\pi}{\mu}(x) \right] \, h(y).
\end{align*}

Combining the transformed left-hand and right-hand sides together gives
%
\begin{align*}
\int \pi( \mathrm{d} x ) \, h(f(x)) 
&= 
\int f_{*} \pi( \mathrm{d} y )  \, h(y)
\\
\int \nu( \mathrm{d} y ) \, 
\left[ \int \mu^{f, \nu}( \mathrm{d} x \mid y ) \, \rnd{\pi}{\mu}(x) \right] \, h(y)
&=
\int \nu( \mathrm{d} y ) \, \left[ \rnd{f_{*}\pi}{\nu}(y) \right] \, h(y).
\end{align*}
%
Because both sides of this equation are $\nu$-informed integrals we have equality if and only if the 
integrands are equal up to $\nu$-null subsets. In particular we have equality for any integrand 
$h : Y \rightarrow \mathbb{R}$ if and only if
%
\begin{equation*}
\rnd{f_{*}\pi}{\nu}(y) 
\overset{\nu}{=}
\int \mu^{f, \nu}( \mathrm{d} x \mid y ) \, \rnd{\pi}{\mu}(x).
\end{equation*}

In theory this gives us an explicit formula for pushforward probability density functions.  
Implementing this result in practice, however, will not straightforward unless we happen 
to have an explicit method to integrate the initial probability density function over each 
level set of $f$.

\section{Simplifications For Product Spaces}

Conditional probability theory can be intimidating in its full generality, but it becomes 
much more straightforward when applied to the natural projection functions of product 
spaces.  Consider, for example, a product space consisting of $I$ component spaces,
%
\begin{equation*}
X^{1:I} = \times_{i = 1}^{I} X_{i},
\end{equation*}
%
where 
%
\begin{equation*}
1 \! : \! I = \{1, \ldots, I\}
\end{equation*}
%
denotes the component indices.

Given any subset of component indices 
%
\begin{equation*}
\mathsf{i} = \{ i_{1}, \ldots, i_{J} \} \subset 1 \! : \! I
\end{equation*}
% 
we can decompose any product variable 
%
\begin{equation*}
x_{1:I} = (x_{1}, \ldots, x_{I})
\end{equation*}
%
into a marginal variable
%
\begin{equation*}
x_{\mathsf{i}} = (x_{i_{1}}, \ldots, x_{i_{J}}) \in \times_{i \in \mathsf{i} } X_{i}
\end{equation*}
%
and a cross sectional variable
%
\begin{equation*}
x_{\setminus \mathsf{i}} 
= ( x_{i} \mid i \notin \mathsf{i} ) 
\in \{ x_{\mathsf{i}} \} \times X^{\setminus \mathsf{i}},
\end{equation*}
%
where
%
\begin{equation*}
\setminus \mathsf{i} = \{ i \mid i \notin \mathsf{i} \}.
\end{equation*}

This decomposition is accompanied by the projection function
%
\begin{equation*}
\varpi_{\mathsf{i}} : X^{1:I} \rightarrow X^{\mathsf{i}}
\end{equation*}
%
with the level sets
%
\begin{equation*}
\varpi_{\mathsf{i}}^{-1}( x_{\mathsf{i}} ) 
= \{ x_{\mathsf{i}} \} \times X^{\setminus \mathsf{i}}.
\end{equation*}

Reference measures on each component space immediately define a product reference 
measure over the product space,
%
\begin{equation*}
\nu^{I} = \times_{i = 1}^{I} \nu_{i}.
\end{equation*}
%
At the same time the product measure
%
\begin{equation*}
\nu^{\mathsf{i}} = \times_{i \in \mathsf{i} } \nu_{i}
\end{equation*}
%
defines a reference measure over the marginal space $X^{\mathsf{i}}$ which also 
happens to be the output space of the projection function $\varpi_{\mathsf{i}}$.
Finally the product of the remaining component measures,
%
\begin{equation*}
\nu^{\setminus \mathsf{i}} = \times_{i \in \setminus \mathsf{i} } \nu_{i},
\end{equation*}
%
defines a reference measure for each cross section 
$\{ x_{\mathsf{i}} \} \times X^{\setminus \mathsf{i}}$, and hence each level 
set of the projection function.  In summary by combining the component 
reference measures together appropriately we can construct reference measures
for the input space, output space, and level sets of the projection function 
$\varpi_{\mathsf{i}}$!

\textbf{Product measure disintegrates into product measure over conditional 
components}
%
\begin{equation*}
(\nu^{I})^{\varpi_{\mathsf{i}}, \nu^{\mathsf{i}}} = \nu^{\setminus \mathsf{i}}.
\end{equation*}

These reference measures allow us to construct conditional probability 
density functions over the cross sections that satisfy the product rule
%
\begin{align*}
\rnd{\pi}{ \nu^{I} }(x_{1:I})
&\overset{ \nu^{I} }{ = }
\rnd{ (\varpi_{\mathsf{i}})_{*} \pi }{\nu^{\mathsf{i}}} (\varpi_{\mathsf{i}}(x_{1:I})) \,
\rnd{ \pi^{\varpi_{\mathsf{i}}} }{ (\nu^{I})^{\varpi_{\mathsf{i}}, \nu^{\mathsf{i}}} } (x_{\setminus \mathsf{i}}  \mid \varpi_{\mathsf{i}}(x_{1:I}))
\\
\rnd{\pi}{ \nu^{I} }(x_{1:I})
&\overset{ \nu^{I} }{ = }
\rnd{ (\varpi_{\mathsf{i}})_{*} \pi }{\nu^{\mathsf{i}}} (\varpi_{\mathsf{i}}(x_{1:I})) \,
\rnd{ \pi^{\varpi_{\mathsf{i}}} }{ \nu^{\setminus \mathsf{i}} } (x_{\setminus \mathsf{i}}  \mid \varpi_{\mathsf{i}}(x_{1:I}))
\\
\rnd{\pi}{ \nu^{I} }(x_{1:I})
&\overset{ \nu^{I} }{ = }
\rnd{ (\varpi_{\mathsf{i}})_{*} \pi }{\nu^{\mathsf{i}}} (x_{\mathsf{i}}) \,
\rnd{ \pi^{\varpi_{\mathsf{i}}} }{ \nu^{\setminus \mathsf{i}} } (x_{\setminus \mathsf{i}}  \mid x_{\mathsf{i}}).
\end{align*}

For example if $X^{1:I} = \mathbb{R}^{6}$ and we take
%
\begin{equation*}
\mathsf{i} = \{ 2, 5, 6 \},
\end{equation*}
%
\begin{equation*}
\setminus \mathsf{i} = \{ 1, 3, 4 \},
\end{equation*}
%
and
%
\begin{equation*}
\nu_{i} = \lambda
\end{equation*}
%
then we have
%
\begin{equation*}
\rnd{\pi}{ \lambda^{6} }( x_{1}, x_{2}, x_{3}, x_{4}, x_{5}, x_{6} )
\overset{ \nu^{I} }{ = }
\rnd{ (\varpi_{\mathsf{i}})_{*} \pi }{ \lambda^{3} } ( x_{2}, x_{5}, x_{6}  ) \,
\rnd{ \pi^{\varpi_{\mathsf{i}}} }{ \lambda^{3} } (x_{1}, x_{3}, x_{4} \mid x_{2}, x_{5}, x_{6} )
\end{equation*}
%
This allows us to for example compute conditional probability densities from the evaluations of 
joint and marginal probability density functions,
%
\begin{equation*}
\rnd{ \pi^{\varpi_{\mathsf{i}}} }{ \lambda^{3} } (x_{1}, x_{3}, x_{4} \mid x_{2}, x_{5}, x_{6} )
\overset{ \nu^{I} }{ = }
\frac{ \rnd{\pi}{ \lambda^{6} }( x_{1}, x_{2}, x_{3}, x_{4}, x_{5}, x_{6} ) }
{ \rnd{ (\varpi_{\mathsf{i}})_{*} \pi }{ \lambda^{3} } ( x_{2}, x_{5}, x_{6} ) }.
\end{equation*}


Similarly
%
\begin{align*}
\rnd{ (\varpi_{\mathsf{i}})_{*}\pi }{ \lambda^{3} }( x_{2}, x_{5}, x_{6}  ) 
&\overset{\nu}{=}
\int (\lambda^{6})^{\varpi_{\mathsf{i}}, \lambda^{3}}( \mathrm{d} x_{1:6} \mid x_{2}, x_{5}, x_{6} ) \, 
\rnd{\pi}{\lambda^{6}}( x_{1}, x_{2}, x_{3}, x_{4}, x_{5}, x_{6} )
\\
&\overset{\nu}{=}
\int \lambda^{3}(x_{1}, x_{3}, x_{4} ) \, 
\rnd{\pi}{\lambda^{6}}( x_{1}, x_{2}, x_{3}, x_{4}, x_{5}, x_{6} )
\\
&\overset{\nu}{=}
\int \mathrm{d} x_{1} \, \mathrm{d} x_{3} \, \mathrm{d} x_{4} \,
\rnd{\pi}{\lambda^{6}}( x_{1}, x_{2}, x_{3}, x_{4}, x_{5}, x_{6} ).
\end{align*}

Given a joint probability density function we can use this result to compute the marginal 
probability density function, assuming we can evaluate the integral analytically, and then 
the product rule to compute the conditional probability density function.  At the same time 
any marginal probability density function paired with a conditional probability density 
function defines a joint probability density function.

Addition example.

\section{Convolutions}

Probability space $(X, \mathcal{X}, \pi)$.

We can construct the second-order power space $X^{2} = X \times X$ with projection 
functions $\varpi_{1}$ and $\varpi_{2}$ that return the first and second identical 
components, respectively.  Any conditional probability distribution over $X^{2}$ 
with respect to $\varpi_{1}$ lifts $\pi$ to a joint probability distribution
%
\begin{equation*}
\pi( \mathrm{d} x_{1}, \mathrm{d} x_{2}) = \pi( \mathrm{d} x_{1} \mid x_{2} ) \, \pi( \mathrm{d} x_{2} ).
\end{equation*}

By construction pushing this joint probability distribution forward along $\varpi_{1}$ 
returns the original probability distribution,
%
\begin{equation*}
(\varpi_{1})_{*} \pi ( \mathrm{d} x_{2}) = \pi( \mathrm{d} x_{2}).
\end{equation*}
%
Pushing it forward along the second projection function $\varpi_{2}$, however, generally 
gives a completely different probability distribution over $X$.

This process of lifting and then projecting to construct new probability distributions 
is known as \textbf{convolution}.

Comparison to additive pushforward.

\end{document}  